\section{The cyclic sum at every conormal depth}
\label{sec:cyclic-sum-conormal}

This section continues the complete cyclic-sum source, especially its
Sections~1--3 and 7--8, and the separately supplied conormal-tower proof.
It retains the original packet, arithmetic unit, tensor variables and
reflection.  The new constructions are the retraction and invariant
dimensions at every depth, the exact scalar first-jet image, its full
conormal comparison, and the finite cyclic Weyl correction with its
complementary map.

\subsection{Original coordinates and the exact collided annihilator}

Let \(k\geq1\), and let the nonempty original monic packet be
\[
 h(s)=\prod_{\nu=1}^{a}(s-\rho_\nu)^{m_\nu},
 \qquad d=\sum_{\nu=1}^{a}m_\nu\geq1.
 \tag{CC.1}
\]
The \(\rho_\nu\) are distinct, and every \(m_\nu\) is the full order
retained by the packet.  In the arithmetic application
\(g=2\xi\), \(v_h=g/h\), and the chosen packet consists of zeros of
\(g\) with those full orders.  Thus \(v_h(\rho_\nu)\ne0\).
The polynomial-algebra statements below apply to the displayed original
roots without any assumption that they lie on the critical line.
Put
\[
\begin{split}
 P&=\mathbb C[s_1,\ldots,s_k],\quad h_i=h(s_i),\quad I=(h_1,\ldots,h_k),
 \quad S=\sum_{i=1}^k s_i,\\
 B_r&=P/I^r,\qquad A_r=B_r^{\mathfrak S_k}\quad(r\geq1),\\
 U_r&=\left[\prod_{i=1}^k v_h(s_i)\right]_{I^r}\in A_r^\times .
\end{split}
\tag{CC.2}
\]
An entire function in (CC.2) means its full finite Taylor class on
each local factor.  These classes are compatible under every quotient
map.  Their invertibility follows from their nonzero constant terms:
in a finite local algebra a class \(u_0+n\), with \(u_0\ne0\) and
\(n^N=0\), has inverse
\(u_0^{-1}\sum_{j=0}^{N-1}(-n/u_0)^j\).
This fixes \(U_r\) without replacing any of its Taylor coefficients.

For an ordered tuple \({\bf t}=(\nu_1,\ldots,\nu_k)\), write
\[
\begin{split}
 \lambda_{\bf t}&=\sum_i\rho_{\nu_i},&
 D_{\bf t}&=\sum_i(m_{\nu_i}-1),&
 M_{\bf t}&=\max_i m_{\nu_i},\\
 L_{{\bf t},r}&=1+D_{\bf t}+(r-1)M_{\bf t},&
 \ell_{\lambda,r}&=\max_{\lambda_{\bf t}=\lambda}L_{{\bf t},r},\\
 \Lambda_k&=\{\lambda_{\bf t}\},&
 \chi_r(T)&=\prod_{\lambda\in\Lambda_k}(T-\lambda)^{\ell_{\lambda,r}},
 \qquad q_r=\deg\chi_r .
\end{split}
\tag{CC.3}
\]
The letter \(T\) is the polynomial indeterminate in \(\chi_r\);
its evaluation in the original algebra is \(T=S\).
Equal sums are grouped before taking the maximum.

\begin{theorem}[Full local depth and the exact scalar contraction]
There are algebra isomorphisms in the original shifted coordinates
\[
 B_r\ \longrightarrow\
 \prod_{{\bf t}\in\{1,\ldots,a\}^k}
 \mathbb C[z_1,\ldots,z_k]/(z_1^{m_{\nu_1}},\ldots,z_k^{m_{\nu_k}})^r,
 \qquad s_i\longmapsto\rho_{\nu_i}+z_i.
 \tag{CC.4}
\]
In the \({\bf t}\)-factor an exact basis consists of the monomials
\[
 z^\alpha,\qquad
 \sum_i\left\lfloor\frac{\alpha_i}{m_{\nu_i}}\right\rfloor\leq r-1.
 \tag{CC.5}
\]
The nilpotency order of \(S-\lambda_{\bf t}=\sum_i z_i\) on that
factor is exactly \(L_{{\bf t},r}\).  Consequently
\[
 I^r\cap\mathbb C[S]=(\chi_r(S)),\qquad
 C_r:=\mathbb C[T]/(\chi_r)
 \ \xrightarrow{\ \alpha_r\ }\ A_r,\quad [F]\longmapsto[F(S)]
 \tag{CC.6}
\]
is an injective unital algebra map.  The original arithmetic coefficient
map on the cyclic source is
\(\eta_r=M_{U_r}\alpha_r:C_r\longrightarrow A_r\).
It is injective and intertwines multiplication by \(T\) and \(S\).
\end{theorem}
\begin{proof}
For \(r=1\), separate-variable Chinese remainder division gives
\(I=\bigcap_{\bf t}J_{\bf t}\), where
\(J_{\bf t}=((s_i-\rho_{\nu_i})^{m_{\nu_i}})_i\).
Two different tuples give comaximal ideals: in a coordinate where
their roots differ the two one-variable powers are coprime, and their
Bezout identity already gives \(1\) in the sum of the two ideals.
For finitely many pairwise comaximal ideals, intersection equals
product.  Therefore
\(I^r=(\prod J_{\bf t})^r=\prod J_{\bf t}^r=\bigcap J_{\bf t}^r\).
The same Bezout identity, or its power expansion, proves that the
powers remain comaximal.  This proves (CC.4).

The local ideal is monomial.  A monomial belongs to its \(r\)-th
power precisely when some nonnegative integers \(b_i\), with
\(\sum b_i=r\), satisfy \(m_{\nu_i}b_i\leq\alpha_i\) for every \(i\).
Such \(b_i\) exist precisely when the sum of the displayed floors is
at least \(r\).  This proves (CC.5), including independence.
Write \(\alpha_i=m_{\nu_i}b_i+c_i\), \(0\leq c_i<m_{\nu_i}\).
The greatest surviving total degree is
\(D_{\bf t}+(r-1)M_{\bf t}\): put every allowed block \(b_i\)
in one coordinate attaining \(M_{\bf t}\), and put
\(c_i=m_{\nu_i}-1\) in every coordinate.
All monomials of larger degree vanish.  In
\((\sum z_i)^{D_{\bf t}+(r-1)M_{\bf t}}\) that chosen monomial
has the nonzero coefficient
\((D_{\bf t}+(r-1)M_{\bf t})!/\prod\alpha_i!\).
Distinct monomials cannot cancel.  The asserted exact order follows.
The annihilator in \(\mathbb C[T]\) of evaluation on this factor is
therefore \(((T-\lambda_{\bf t})^{L_{{\bf t},r}})\).
Intersecting these ideals takes the maximum at each repeated
\(\lambda\), proving (CC.6).  The class of \(F(S)\) is invariant,
and multiplication by the unit \(U_r\) preserves injectivity and
the stated intertwining.
\end{proof}

\begin{proposition}[Degree, scalar powers and all endpoints]
For the nonempty packet,
\[
\begin{split}
 q_r&\geq k(d-1)+(r-1)d+1,\\
 \chi_r&\mid\chi_1^r,\qquad
 \chi_r=\chi_1^r
 \ \Longleftrightarrow\
 (r=1)\ \text{or}\ (k=1)\ \text{or}\ (m_\nu=1\ \forall\nu).
\end{split}
\tag{CC.7}
\]
These are assertions about the complete collided polynomials.
\end{proposition}
\begin{proof}
Repeated monic division in the unchanged variables gives the unique
finite expansion
\[
 F=\sum_{{\bf b}\in\mathbb N^k,\ 0\leq c_i<d}
 a_{{\bf b},{\bf c}}\prod_i h_i^{b_i}s_i^{c_i}.
 \tag{CC.8}
\]
The leading monomials have the pairwise distinct exponent vectors
\(d{\bf b}+{\bf c}\), so a finite dependence is impossible.
Equivalently \(P\) is free over \(\mathbb C[h_1,\ldots,h_k]\)
with the displayed remainder basis.  Membership in \(I^r\)
is exactly the vanishing of the coefficients with \(\sum b_i<r\).
It follows that the highest homogeneous part of every element of
\(I^r\) belongs to \((s_1^d,\ldots,s_k^d)^r\).
Indeed each highest-degree basis term has that property, and its
lower terms have smaller total degree.

Put \(D=k(d-1)+(r-1)d\).
The monomial with exponents \((rd-1,d-1,\ldots,d-1)\) survives the
corresponding monomial ideal and has degree \(D\).
For each \(0\leq j\leq D\), choose a divisor of this monomial of
degree \(j\).  Its floor sum is still at most \(r-1\), and its
coefficient in \(S^j\) is a nonzero multinomial coefficient.
Thus the leading homogeneous part \(S^j\) of a monic polynomial
of degree \(j\) in \(S\) cannot belong to that monomial ideal.
Such a polynomial cannot lie in \(I^r\).  Apply this to \(\chi_r(S)\)
to obtain the lower bound.

For every tuple, \(M_{\bf t}\leq1+D_{\bf t}=L_{{\bf t},1}\), so
\(L_{{\bf t},r}\leq rL_{{\bf t},1}\).
Taking maxima proves divisibility.  Each of the three stated
equality cases follows directly from (CC.3).
Conversely suppose \(r,k>1\), and put \(M=\max_\nu m_\nu>1\).
Choose a root \(\rho\) of order \(M\) and its pure \(k\)-tuple.
At \(\lambda=k\rho\), every tuple has
\(D_{\bf t}\leq k(M-1)\) and \(M_{\bf t}\leq M\);
the pure tuple attains both bounds.  Therefore
\[
 r\ell_{\lambda,1}-\ell_{\lambda,r}
 =(r-1)(k-1)(M-1)>0.
 \tag{CC.9}
\]
This proves strict divisibility even when other tuples collide there.
\end{proof}

For one root of order \(m\), (CC.3) reads
\(\chi_r(T)=(T-k\rho)^{1+k(m-1)+(r-1)m}\).
For \(k=1\), \(\chi_r=h^r\) and \(\alpha_r\) is an isomorphism.
For \(d=1\), \(\chi_r=(T-k\rho)^r\), but the invariant algebra
can have a complement to the cyclic image: for \(k\geq2,r=3\),
the two symmetric degree-two polynomials
\(\sum z_i^2\) and \(\sum_{i<j}z_iz_j\) are independent.
The empty packet \(h=1\) is treated separately: \(I=P\),
\(B_r=A_r=C_r=0\), \(\chi_r=1\), and every displayed finite linear
map is the zero map.  No formula involving a chosen root or
\(q_r^{-1}\) is applied in that case.  The arithmetic source itself
need not be zero merely because its finite packet projection is zero.

The attaining tuple can change with depth.  As a declared algebraic
example, take roots \(1/2+it\), \(t=0,1,2,3\), of orders \(6,4,4,1\),
and \(k=2\).  At \(\lambda=1+3i\) the only unordered pairs are
\((0,3)\) and \((1,2)\); they give respectively \(6r\) and \(4r+3\).
Thus \(\ell_{\lambda,r}=\max(6r,4r+3)\).
This example asserts no location or multiplicity of an arithmetic zero.

\subsection{Unscaled invariant bases and a full-unit retraction}

For an occupation vector \({\bf n}=(n_1,\ldots,n_a)\) with
\(\sum n_\nu=k\), put \(\lambda_{\bf n}=\sum n_\nu\rho_\nu\).
In the factor for a fixed ordered tuple of that occupation, its
stabilizer is \(\prod_\nu\mathfrak S_{n_\nu}\).
For each orbit of the surviving monomials (CC.5), take the sum of
the distinct monomials in that orbit, with coefficient one.
These are the original unscaled invariant basis elements.
Their exact number is
\[
 a_{{\bf n},r}
 =\sum_{p=0}^{r-1}
 [t^p x_1^{n_1}\cdots x_a^{n_a}]
 \prod_{\nu=1}^a\prod_{j=0}^{r-1}
 (1-x_\nu t^j)^{-m_\nu},
 \qquad
 a_{\lambda,r}=\sum_{\lambda_{\bf n}=\lambda}a_{{\bf n},r}.
 \tag{CC.10}
\]
Coefficient extraction here is finite formal algebra.  Each variable
in a group of order \(m_\nu\) has a pair of labels
\((j,c)\), \(j\geq0,\ 0\leq c<m_\nu\), representing the original
exponent \(m_\nu j+c\).  An orbit is a multiset of these labels,
and the surviving condition is that their \(j\)'s sum to at most
\(r-1\).  The geometric series for each label proves (CC.10).
An invariant in the full product (CC.4) is determined by its entry
at one ordered tuple per occupation; the other entries are its
variable-permuted translates.  The stabilizer condition is precisely
the one just imposed.  This proves both the claimed basis and
\[
\begin{split}
 \dim B_r&=d^k\binom{r+k-1}{k},\\
 \dim A_r&=\sum_\lambda a_{\lambda,r}
 =\sum_{p=0}^{r-1}[t^p x^k]
       \prod_{j=0}^{r-1}(1-xt^j)^{-d},\\
 a_{{\bf n},1}&=\prod_\nu\binom{n_\nu+m_\nu-1}{m_\nu-1}.
\end{split}
\tag{CC.11}
\]
For the first equality use (CC.8): there are
\(\binom{r+k-1}{k}\) vectors \({\bf b}\) with \(\sum b_i<r\),
and \(d^k\) possible remainders.  For the second alternative expression,
permutation acts on that same basis by permuting its \(k\) pairs
\((b_i,c_i)\); its unscaled orbit sums give the stated series.

\begin{theorem}[Constructive retraction at every depth]
There is an explicit \(\mathbb C[T]\)-linear retraction
\(\Pi_r:A_r\longrightarrow C_r\) of \(\eta_r\), with the complete
arithmetic unit retained.  Put \(K_r=\ker\Pi_r\) and
\(Q_r=1-\eta_r\Pi_r\).  Then
\[
\begin{split}
 C_r\oplus K_r&\longrightarrow A_r,& (c,v)&\longmapsto\eta_r(c)+v,\\
 A_r&\longrightarrow C_r\oplus K_r,& w&\longmapsto(\Pi_rw,Q_rw)
\end{split}
\tag{CC.12}
\]
are inverse \(\mathbb C[T]\)-module isomorphisms.  In particular
\(A_r/\eta_rC_r\longrightarrow K_r,\ [w]\longmapsto Q_rw\)
is an isomorphism, inverse to inclusion followed by quotient.
\end{theorem}
\begin{proof}
For each \(\lambda\), use the exact CRT idempotent \(e_\lambda(T)\)
in \(C_r\), and let \(A_{\lambda,r}=e_\lambda(S)A_r\).
Write \(N=S-\lambda\) on this summand.
Choose an ordered tuple attaining \(\ell=\ell_{\lambda,r}\).
Choose a coordinate of maximal multiplicity in it, put all \(r-1\)
allowed blocks in that coordinate, and put every remainder at its
maximum \(m_{\nu_i}-1\).  Call the resulting exponent \(\alpha^*\).
Thus \(|\alpha^*|=\ell-1\).
Let \(\theta:A_{\lambda,r}\to\mathbb C\) extract the coefficient
of \(z^{\alpha^*}\) in that specific local factor, in the basis (CC.5).
Define, with \(X=T-\lambda\),
\[
\begin{split}
 \Pi_\lambda^0(v)&=\sum_{j=0}^{\ell-1}
       \theta(N^{\ell-1-j}v)X^j\pmod{X^\ell},\\
 w_\lambda&=e_\lambda(S)U_r,\qquad
 p_\lambda(X)=\Pi_\lambda^0(w_\lambda),\\
 \Pi_\lambda(v)&=p_\lambda(X)^{-1}\Pi_\lambda^0(v).
\end{split}
\tag{CC.13}
\]
The equality \(N^\ell=0\) gives
\(\Pi_\lambda^0(Nv)=X\Pi_\lambda^0(v)\), coefficient by coefficient:
the constant term on the left is zero and the other coefficients
shift by one.  Also
\[
 p_\lambda(0)=
 \left(\prod_i v_h(\rho_{\nu_i})\right)
 \frac{(\ell-1)!}{\prod_i\alpha_i^*!}\ne0.
 \tag{CC.14}
\]
To verify (CC.14), the selected term in \(N^{\ell-1}\) has that
multinomial coefficient; multiplying by a positive-degree term of
the unit would exceed the greatest surviving total degree.
The idempotent \(e_\lambda(S)\) is exactly one on this local factor.
Set \(t=p_\lambda/p_\lambda(0)-1\).  Its exact inverse is
\(p_\lambda(0)^{-1}\sum_{j=0}^{\ell-1}(-t)^j\).
This retains every coefficient of \(p_\lambda\).
The intertwining now gives
\(\Pi_\lambda(e_\lambda(S)U_rF(S))=F(\lambda+X)\bmod X^\ell\).
Take the product of these maps and use scalar CRT to define \(\Pi_r\).
It is the asserted retraction.  Its kernel is \(S\)-stable.
The identities \(Q_r\eta_r=0\), \(\Pi_rQ_r=0\), and
\(Q_r|_{K_r}=1\) prove every map and inverse in (CC.12).
\end{proof}

For any polynomial or entire \(f\), multiplication by \(f(S)\) has
on \(A_{\lambda,r}\) the form \(f(\lambda)1\) plus a polynomial
without constant term in the nilpotent \(N\).  Powers of a nilpotent
map have trace zero, as is seen in a basis adapted to its kernel
filtration.  The same statement holds on the two invariant summands
in (CC.12).  Thus
\[
\begin{split}
 \operatorname{Tr}_{A_r}f(S)&=\sum_\lambda a_{\lambda,r}f(\lambda),\\
 \operatorname{Tr}_{C_r}f(T)&=\sum_\lambda\ell_{\lambda,r}f(\lambda),\\
 \operatorname{Tr}_{K_r}f(S)&=
       \sum_\lambda(a_{\lambda,r}-\ell_{\lambda,r})f(\lambda).
\end{split}
\tag{CC.15}
\]
This also proves \(a_{\lambda,r}\geq\ell_{\lambda,r}\).
All nilpotent coordinates remain in (CC.4), (CC.10) and (CC.13);
the trace formula follows from those actual operators.

\subsection{The derivative tower, scalar first jets, and the unit obstruction}

Denote the quotient maps by
\(\pi_r:B_{r+1}\to B_r\) and \(\rho_r:C_{r+1}\to C_r\).
The original derivative is
\[
 \delta_r:B_{r+1}\longrightarrow B_r,\quad
 [F]\longmapsto\left[\frac1k\sum_i\partial_{s_i}F\right],
 \qquad
 \partial_r:C_{r+1}\longrightarrow C_r,\quad [F]\longmapsto[F'].
 \tag{CC.16}
\]
Each is a complex linear map.  Differentiating a product of \(r+1\)
elements of \(I\) leaves at least \(r\) such factors, so
\(\partial_S I^{r+1}\subset I^r\).
Formula (CC.3) gives
\(\ell_{\lambda,r+1}\geq\ell_{\lambda,r}+1\).
Thus the derivative of a multiple of \(\chi_{r+1}\) is a multiple
of \(\chi_r\), proving the second descent as well.  On representatives,
\[
\begin{split}
 \delta_r\alpha_{r+1}&=\alpha_r\partial_r,&
 \pi_r\alpha_{r+1}&=\alpha_r\rho_r,\\
 \delta_r(xy)&=\pi_r(x)\delta_r(y)+\pi_r(y)\delta_r(x),&
 \pi_r\delta_{r+1}&=\delta_r\pi_{r+1}.
\end{split}
\tag{CC.17}
\]
The last equality is a map \(B_{r+2}\to B_r\).
The symmetric subspaces are preserved because the derivative averages
all \(k\) coordinates with the literal factor \(1/k\).
Every \(\delta_r\) is onto: the invertible coordinates
\(S,z_1,\ldots,z_{k-1}\), where \(z_i=s_i-S/k\), have
\(\partial_S S=1\) and \(\partial_Sz_i=0\).  Integrating powers
of \(S\) gives a polynomial primitive of any polynomial representative.
Averaging such a primitive proves surjectivity on the invariant
subspaces too.

The same monic expansion (CC.8) identifies
\[
 I^j/I^{j+1}\cong
 (P/I)[\zeta_1,\ldots,\zeta_k]_j,\quad
 \zeta_i\longmapsto[h_i],
 \qquad
 \operatorname{gr}\delta
 =\frac1k\sum_i[h_i']\,\partial_{\zeta_i}.
 \tag{CC.18}
\]
Indeed differentiating \(h^{\bf b}r_{\bf b}\) in degree \(j=|{\bf b}|\)
gives \(b_i h^{{\bf b}-e_i}h_i'r_{\bf b}/k\) in degree \(j-1\);
differentiating \(r_{\bf b}\) stays in degree \(j\).
Both the integer \(b_i\) and the original polynomial \(h_i'\) survive.

Introduce a formal square-zero parameter \(\epsilon\), with its
ordinary coefficient algebra meaning.  The two maps
\[
\begin{split}
 \Phi_r:B_{r+1}&\longrightarrow B_r[\epsilon]/(\epsilon^2),
 &b&\longmapsto\pi_rb+\epsilon\delta_rb,\\
 \phi_r:C_{r+1}&\longrightarrow C_r[\epsilon]/(\epsilon^2),
 &F&\longmapsto\rho_rF+\epsilon\partial_rF
\end{split}
\tag{CC.19}
\]
are unital algebra homomorphisms by (CC.17) and its one-variable
product rule.  Extending \(\alpha_r\) coefficientwise to
\(\alpha_r^\epsilon\) gives the exact square
\(\Phi_r\alpha_{r+1}=\alpha_r^\epsilon\phi_r\).
The coefficient projections of these algebra maps are precisely
the quotient and derivative maps in (CC.16).

Put \(R(T)=\prod_{\lambda\in\Lambda_k}(T-\lambda)\) and
\(\psi_r=\chi_rR\).  Then
\[
\begin{split}
 \ker\phi_r&=(\psi_r)/(\chi_{r+1}),&
 \dim\operatorname{im}\phi_r&=q_r+|\Lambda_k|,\\
 \ker\rho_r&=(\chi_r)/(\chi_{r+1})
       \cong\mathbb C[T]/(\chi_{r+1}/\chi_r),&
 \partial_r(\ker\rho_r)&=\chi_r'C_r .
\end{split}
\tag{CC.20}
\]
For the first assertion, \(F\) and \(F'\) vanish modulo
\((T-\lambda)^{\ell_{\lambda,r}}\) exactly when \(F\) has order
at least \(\ell_{\lambda,r}+1\) there.  The order implication in
(CC.16) gives \(\psi_r\mid\chi_{r+1}\), so the asserted kernel is
well defined and has dimension \(q_{r+1}-q_r-|\Lambda_k|\).
This proves the image dimension and the injective factorization
\(\mathbb C[T]/(\psi_r)\hookrightarrow C_r[\epsilon]/(\epsilon^2)\).
Multiplication by \(\chi_r\) gives the displayed isomorphism onto
\(\ker\rho_r\); its kernel is exactly \((\chi_{r+1}/\chi_r)\).
Differentiating \(\chi_rF\) modulo \(\chi_r\) gives \(\chi_r'F\).
At each \(\lambda\), multiplication by \(\chi_r'\) maps the input
constant to a nonzero multiple of
\((T-\lambda)^{\ell_{\lambda,r}-1}\).
Its image has one dimension per \(\lambda\), proving all image
and kernel claims even at repeated collided roots.

Write
\[
 \beta_r=U_r^{-1}\delta_rU_{r+1}\in A_r,\qquad
 \Phi_r(U_{r+1})=U_r(1+\epsilon\beta_r),\qquad
 \Phi_r(U_{r+1}^{-1})=U_r^{-1}(1-\epsilon\beta_r).
 \tag{CC.21}
\]
The first formula is the definition of the full logarithmic derivative
class, the second follows directly from (CC.19), and the third follows
by multiplication and \(\epsilon^2=0\).
In particular,
\[
 \delta_r\eta_{r+1}
   =\eta_r\partial_r+
      M_{\delta_rU_{r+1}}\alpha_r\rho_r.
 \tag{CC.22}
\]
This formula retains the full derivative of the arithmetic unit.

There is a complete decomposition of its last term using (CC.12):
\[
\begin{split}
 b_r&=\Pi_r(\delta_rU_{r+1})\in C_r,&
 n_r&=Q_r(\delta_rU_{r+1})\in K_r,\\
 \mathcal N_r:C_r&\longrightarrow K_r,&
 \mathcal N_r(F)&=\alpha_r(F)n_r,\\
 \delta_r\eta_{r+1}
   &=\eta_r(\partial_r+M_{b_r}\rho_r)+\mathcal N_r\rho_r .
\end{split}
\tag{CC.23}
\]
Here \(K_r\) is a \(\mathbb C[T]\)-submodule, so multiplication
by \(\alpha_r(F)\) preserves it.  Both \(\Pi_r\) and \(Q_r\)
intertwine this multiplication; applying them to the last term
of (CC.22) proves (CC.23).
To state the obstruction independently of the chosen retraction, let
\(\mathfrak q_r:A_r\to A_r/\eta_rC_r\) be the module quotient and put
\[
 \Omega_r=\mathfrak q_r M_{\delta_rU_{r+1}}\alpha_r:
 C_r\longrightarrow A_r/\eta_rC_r.
 \quad
 \mathfrak q_r\delta_r\eta_{r+1}=\Omega_r\rho_r,
 \qquad
 \Omega_r=0\ \Longleftrightarrow\ \beta_r\in\alpha_r(C_r).
 \tag{CC.24}
\]
To prove the equivalence, evaluate \(\Omega_r\) at \(1\).
Its vanishing says \(\delta_rU_{r+1}=U_r\alpha_r(\gamma_r)\)
for a unique \(\gamma_r\in C_r\), precisely the condition on \(\beta_r\).
Conversely that equality makes every value of \(\Omega_r\) zero
by closure of the subalgebra \(\alpha_r(C_r)\).
Under (CC.12), \(\Omega_r\) corresponds exactly to \(\mathcal N_r\).
When it vanishes, \(b_r=\gamma_r\), and the exact pulled-back
derivative is \(\partial_r+M_{\gamma_r}\rho_r\).
The formulas construct both the obstruction and its strongest
scalar derivative correspondence without deleting the complementary
directions.
